\section{微积分基本定理}
\label{sec:ru-fundamental-theorem}

求导记录局部变化，积分记录一段范围内的累计。两者能否相互恢复，要分两个方向回答：对累计函数求导能否得到被积函数；对导数积分能否得到原函数的增量。

\subsection{累计函数的变化率}
\label{sec:ru-accumulation-derivative}

设$f$在$[a,b]$上可积，固定积分起点$a$，令
\[
 F(x)=\int_a^xf(t)\,\mathrm dt,\qquad a\leq x\leq b,
\]
称为$f$的累计函数，$t$是积分变量，$x$是$F$的自变量。若$|f|\leq M$，由区间可加性与绝对值估计，$|F(y)-F(x)|=\bigl|\int_x^yf\bigr|\leq M|y-x|$，所以累计函数总是Lipschitz连续的。
差商$(F(x+h)-F(x))/h=\frac1h\int_x^{x+h}f$恰是$f$在短区间上的平均值。若$f$在该点连续，短区间上的平均值应趋于该点的函数值。
\begin{theorem}[累计函数的导数]
\label{thm:ru-accumulation-derivative}
设$f$在$[a,b]$上可积，$F(x)=\int_a^xf$。若$f$在$c\in(a,b)$处连续，则$F'(c)=f(c)$；在端点处连续时，相应的单侧导数等于端点值。特别地，若$f$在$[a,b]$上连续，则$F'=f$处处成立。
\end{theorem}
\begin{proof}
对充分小的$h\neq0$，
\[
 \Bigl|\frac{F(c+h)-F(c)}h-f(c)\Bigr|=\Bigl|\frac1h\int_c^{c+h}\bigl(f(t)-f(c)\bigr)\,\mathrm dt\Bigr|\leq\sup_{|t-c|\leq|h|}|f(t)-f(c)|,
\]
由$f$在$c$处连续，右侧随$h\to0$趋于零。
\end{proof}
这里有三种强弱不同的结论：累计函数总是连续的；它在被积函数的连续点可微，导数等于函数值；被积函数处处连续时，它才处处可微。连续点的条件不能省略。
在$[-1,1]$上取阶跃函数$f=\mathbb I_{[0,1]}$，它只有一个间断点，因而可积，累计函数为$F(x)=\int_{-1}^xf=\max\{x,0\}$。在$x\neq0$处$F'=f$，在$x=0$处左差商恒为$0$、右差商恒为$1$，$F'(0)$不存在。
在间断点处，导数也可能存在却不等于函数值：把$f$取为只在$0$处等于$1$、其余处为$0$的函数，它的累计函数恒为零，$F'(0)=0\neq f(0)$。

\subsection{由原函数计算积分}
\label{sec:ru-newton-leibniz}

上一小节从积分出发得到导数。反过来，若已经知道一个函数的导数是$f$，能否用它计算$f$的积分？
设$f$定义在区间$I$上，若可微函数$G$在$I$上满足$G'=f$，就称$G$为$f$的\term{原函数}（antiderivative）。两个原函数的差的导数为零，由中值定理，它们在同一区间上只相差一个常数。
\begin{theorem}[Newton--Leibniz公式]
\label{thm:ru-newton-leibniz}
设$f$在$[a,b]$上连续，$G$是它在$[a,b]$上的任一原函数，则
\[
 \int_a^bf(x)\,\mathrm dx=G(b)-G(a)=:\bigl[G(x)\bigr]_a^b.
\]
\end{theorem}
\begin{proof}
由定理~\ref{thm:ru-accumulation-derivative}，累计函数$F(x)=\int_a^xf$也是$f$的原函数，所以$G=F+C$。于是$G(b)-G(a)=F(b)-F(a)=\int_a^bf$。
\end{proof}
这个公式把求定积分化为求原函数：例如$(x^3/3)'=x^2$，所以$\int_0^1x^2\,\mathrm dx=1/3$，与图~\ref{fig:ru-riemann-sums}中上下和夹逼的数值一致。

被积函数连续只是一个方便的充分条件。若$G$在$[a,b]$上连续、在$(a,b)$内可导，并且$G'$在$[a,b]$上黎曼可积（端点处任取有限值），则同样有$\int_a^bG'=G(b)-G(a)$。
证明直接用中值定理：对任意分割，在每段上选$\xi_i$使$G(x_i)-G(x_{i-1})=G'(\xi_i)\Delta x_i$，于是$G(b)-G(a)=\sum_iG'(\xi_i)\Delta x_i$恰好是一个黎曼和，令网格尺度趋零，它趋于$\int_a^bG'$。

三个概念需要分开：有原函数、黎曼可积、原函数能用基本函数写出。阶跃函数$\mathbb I_{[0,1]}$在$[-1,1]$上可积，却没有原函数：若$G'=f$，则$G$在$[-1,0]$上为常数、在$[0,1]$上为$x$加常数，在$0$处左导数为$0$、右导数为$1$，与$G'(0)=1$矛盾。
连续函数总有原函数，即它的累计函数；但原函数未必能用基本函数经有限次运算写出。$e^{-x^2}$就是这样的例子（其证明超出本书范围），它的原函数通常直接写成累计函数$\int_0^xe^{-t^2}\,\mathrm dt$。
导数不黎曼可积时能否由导数恢复函数，涉及绝对连续性，放在第~\ref{chap:measurable-functions}章讨论。
